\subsection{紧李代数的嘉当子代数}

引理 1. 设 G 是李群，K 是 G 的紧子群，F 是 L(G) 上的不变双线性型。设 \(x, y \in \mathrm{L}(\mathrm{G})\) 。则存在 K 的元素 k，使得 \(\mathrm{F}(u, [(\mathrm{Ad}k)(x), y]) = 0\) 对一切 \(u \in \mathrm{L}(\mathrm{K})\) 成立。

从 K 到 R 的函数 \(v \mapsto \mathrm{F}((\operatorname{Ad} v)(x), y)\) 连续，因而在某点 \(k \in K\) 处取得极小值。设 \(u \in \mathrm{L}(K)\) ，令

\[
h (t) = \mathrm{F} ((\mathrm{Ad} \exp (t u). k) (x), y), \quad t \in \mathbf {R}.
\]

我们有 \(h(t) \geq h(0)\) （对一切 \(t\) ）；此外，由第 III 章 §3 第 12 节命题 44，

\[
\frac {d h}{d t} (0) = \mathrm{F} ([ u, (\mathrm{Ad} k) (x) ], y) = \mathrm{F} (u, [ (\mathrm{Ad} k) (x), y ]),
\]

引理得证（《实变函数》，第 I 章，§1，第 7 节，命题 7）。

定理 1. 设 g 是紧李代数。g 的嘉当子代数（第 VII 章，§2，第 1 节，定义 1）就是它的极大交换子代数；特别地，g 是它的诸嘉当子代数之并。群 \(\operatorname{Int}(\mathfrak{g})\) 传递地作用在 g 的嘉当子代数所成的集上。

由于 g 是约化的，它的嘉当子代数是交换的（第 VII 章，§2，第 4 节，定理 2 的推论 3）。反之，设 t 是 g 的一个交换子代数。由 §1 第 3 节命题 1，对一切 \(x \in t\) ，ad x 是半单的；由第 VII 章，§2，第 3 节，命题 10，存在 g 的包含 t 的嘉当子代数。这就证明了定理的第一个断言。

现设 \(\mathfrak{t}\) 与 \(\mathfrak{t}'\) 是 \(\mathfrak{g}\) 的两个嘉当子代数。我们证明存在 \(u \in \operatorname{Int}(\mathfrak{g})\) 使得 \(u(\mathfrak{t}) = \mathfrak{t}'\) 。由 §1 第 3 节命题 1，可以假定 \(\mathfrak{g}\) 具有 L(G) 的形式，其中 G 是连通紧李群，并且可以在 g 上选取一个分离的不变对称双线性型 F。设 x（相应地，\(x'\) ）是 g 的正则元，使得 \(\mathfrak{t} = \mathfrak{g}^{0}(x)\) （相应地，\(\mathfrak{t}' = \mathfrak{g}^{0}(x')\) ）（第 VII 章，§3，第 3 节，定理 2）。对 K = G 应用引理 1，可见存在 \(k \in G\) 使得 \([(\operatorname{Ad} k)(x), x']\) 关于 F 与 g 正交，从而为零；于是 \((\operatorname{Ad} k)(x) \in \mathfrak{g}^{0}(x') = \mathfrak{t}'\) ，故 \(\mathfrak{g}^{0}((\operatorname{Ad} k)(x)) = \mathfrak{t}'\) ，因为 \((\operatorname{Ad} k)(x)\) 正则。由此得出 \((\operatorname{Ad} k)(\mathfrak{t}) = \mathfrak{t}'\) ，定理得证。

推论. 设 t 与 \(t'\) 是 g 的嘉当子代数，a 是 t 的子集，u 是把 a 变到 \(t'\) 中的 g 的自同构。则存在 \(\text{Int}(\mathfrak{g})\) 的元素 v，使得 \(u \circ v\) 把 t 变到 \(t'\) ，且在 a 上与 u 相合。

令 \(\mathrm{G} = \operatorname{Int}(\mathfrak{g})\) ，并考虑固定子群 \(Z_{\mathrm{G}}(\mathfrak{a})\) ，即 \(\mathfrak{a}\) 在 \(G\) 中的固定子群；这是 \(G\) 的一个李子群，其李代数 \(\mathfrak{z}_{\mathfrak{g}}(\mathfrak{a})\) 由 \(\mathfrak{g}\) 中与 \(\mathfrak{a}\) 的每个元素交换的元素组成（第 III 章，§9，第 3 节，命题 7）。于是 \(t\) 与 \(u^{-1}(t')\) 是紧李代数 \(\mathfrak{z}_{\mathfrak{g}}(\mathfrak{a})\) 的两个嘉当子代数。由定理 1，存在元素 \(v\) 属于 \(Z_{\mathrm{G}}(\mathfrak{a})\) ，使得 \(v(t) = u^{-1}(t')\) ；任何这样的元素都具有所要的性质。

